\documentclass[twoside, kulak]{kulakreport} % options: kulak (default), kuleuven-brugge of kul

\usepackage[dutch]{babel}

\title{Verslag Practicum}
\subtitle{Ondertitel van het document}
\author{Daria Dehtiar, Adelina Vozianu}
\emailaddress{adelina.vozianu@student.kuleuven.be, daria.dehtiar@student.kuleuven.be}
\date{Academiejaar 2026 -- 2027}
\faculty{Groep Wetenschap \& Technologie Kulak}
\group{Numerieke Wiskunde}
\institute{Informatica}
\address{
	\textbf{KU Leuven Kulak} \\
	\textbf{Groep Wetenschap \& Technologie} \\
	Etienne Sabbelaan 53, 8500 Kortrijk \\
	Tel.\ +32 12 34 56 78 \\
	\href{mailto:\theemailaddress}{\texttt{\theemailaddress}}
}

\begin{document}
	
	\titlepage
	
	\tableofcontents
	
	\chapter*{Inleiding}
	Inleidende tekst.
	
	\chapter{Numerieke technieken om een stelsel op te lossen}
	\section{Inverteren}
	% MOET AANSLUITENDE TEXT IZIJN
	\subsection{Opdracht 1} % zal later weg doen, wanneer we het aaneensluitend tekst maken
	De \texttt{inv()} commando inverteert een gegeven vierkante matrix A. Om geen numerieke nauwkeurigheid te verliezen, is het best om de \texttt{backslash-operator $\backslash$} te benutten. 
	Dit omdat \texttt{inv()} commando 
	
	%Dit omdat de \texttt{backslash-operator} of ook gekend als de \texttt{mldivide()} commando
	
	\chapter*{Besluit}
	Afsluitende tekst.
	
	\bibliographystyle{unsrt}
	\bibliography{bibliografieNW}
	
\end{document}